% Imported from topics/video-generation/seedance-2.5/evidence.md
\chapter{公开能力与证据边界}\label{seedance-m00}
\moduleinfo{M00}{公开能力梳理；沿用来源仓库既有记录}{\hyperref[src-seedance-2.5-release]{\texttt{seedance-2.5-release}}、\hyperref[src-seedance-2.5-product]{\texttt{seedance-2.5-product}}、\hyperref[src-seedance-1.0]{\texttt{seedance-1.0}}、\hyperref[src-seedance-1.5-pro]{\texttt{seedance-1.5-pro}}、\hyperref[src-seedance-2.0]{\texttt{seedance-2.0}}}

本章集中保存前文使用的公开能力依据与边界，供对照总图和候选假设时查询。

先不急着猜 Seedance 2.5 有多少层。看它要解决的问题：同一个人物要出现在后续镜头中，已经发生的动作不能忘掉，声音要跟画面对应，局部出错后还要能修改。只生成一段看起来不错的视频，和把这些要求同时做出来，中间还隔着表示、条件和训练的问题。下面先把公开能力翻译成这些具体问题。

\section{时长增长带来的是上下文问题}\label{seedance-m00-ux65f6ux957fux589eux957fux5e26ux6765ux7684ux662fux4e0aux4e0bux6587ux95eeux9898}

Seedance 2.5 于 2026 年 7 月 31 日正式发布，单次生成时长达到 30 秒，并支持在已有视频上延长。更长的生成窗口让模型有机会在一次输出中安排铺垫、推进和收尾，也增加了保持角色、场景与动作状态的难度。

比如一个人先拿起钥匙，再走到门前。后续镜头要保持的就不只是这张脸，还包括钥匙已经到了手中这一状态。若只比较首尾两张图的相似程度，钥匙在中间消失了也可能没有被发现。所以长叙事中的一致性，需要把外观和事件状态一起考虑。

这里有个不能省略的区别：联合长上下文、分块续写、层次生成和短镜头编排，都可能做出较长成片。产品支持长叙事，只说明最终能力，不足以唯一确定内部采用哪一条路线。前文因而分别讨论长单镜头、多镜头和交互续写，各自检查它需要保持什么，不能只拿总秒数作结论。

\section{更多参考意味着更复杂的绑定}\label{seedance-m00-ux66f4ux591aux53c2ux8003ux610fux5473ux7740ux66f4ux590dux6742ux7684ux7ed1ux5b9a}

单次参考上限为 30 张图片、10 段视频和 10 段音频。这个容量让多角色、多场景和复杂声画要求成为可能，但有效使用这些素材的关键，在于明确每份参考的作用。

具体来说，图片可能只负责人物外观，视频只负责动作与相机，音频只负责音色或节奏。原素材同时还有背景、服装、音乐等信息，它们未必都是目标条件。于是参考问题要多写一步：这份素材中的哪些信息可用，哪些信息需要抑制？只把文件按顺序送进去，并没有回答这个问题。

另外，文件数与条件长度不是一回事。一个视频编码后就可能产生很多 token，继续增加参考，会同时增加编码、交互和显存开销。因此还要决定如何压缩、检索、分配优先级，以及每段特征属于哪个角色。这些关系确定之后，参考上限才有可计算的含义。

\section{编辑能力改变了失败处理方式}\label{seedance-m00-ux7f16ux8f91ux80fdux529bux6539ux53d8ux4e86ux5931ux8d25ux5904ux7406ux65b9ux5f0f}

再看时间戳编辑。假设只有第二到第四秒的服装不对，已经通过的其它画面应当尽量保留。此时任务从重新生成整段，变成了在指定时间、指定对象上改变指定属性。时间戳给出了范围，但怎样让变化不传播到其它部分，还需要另行建模。

这里的“其它不变”还得继续分解。只改颜色时，可以要求外区直接复用原像素；若改成另一个视角，原像素的位置本来就会改变，此时应保持的是身份和空间关系。因此时间区间、目标区域与保持条件要一起给出，不能对所有编辑都套一个像素相等的要求。

白模则先给出布局、姿态、轨迹和相机调度。它可以渲染成帧，也可以导出深度、姿态等结构条件；所以看到“支持白模”，还不能断言主干直接读取三维网格。绿幕任务同样需要拆开：前景怎样保留，背景怎样替换，边缘和光照怎样融合，各自对应不同处理。

\section{音视频联合生成要求更严格的数据}\label{seedance-m00-ux97f3ux89c6ux9891ux8054ux5408ux751fux6210ux8981ux6c42ux66f4ux4e25ux683cux7684ux6570ux636e}

2.5 延续的是统一多模态音视频联合生成方向。这里所说的联合，至少要解释物理时间上的对应关系：嘴形与对白、落地与撞击、手势与演奏节拍，为什么应在某个时刻共同发生？画面和声音主题相近，还没有解决这种同步问题。

输入包含文本、图片、视频和音频，是否就需要四个生成分支？不一定。输入先由编码器变成条件，输出再由相应分支预测，两个数量没有简单的一一对应。双分支是已知可行方向，但 codec、跨模态模块、层数和权重，仍要按每个具体实现分析，不能由产品输入种类反推出来。

\section{怎样把这些能力转成自己的方案}\label{seedance-m00-ux600eux6837ux628aux8fd9ux4e9bux80fdux529bux8f6cux6210ux81eaux5df1ux7684ux65b9ux6848}

从实施角度，可以先选择拿得到的基座，跑通短视频生成，再加入少量图片参考。身份条件确实起作用之后，才继续研究局部编辑和多镜头；原生长上下文与联合音频的数据、算力另外计算。这样安排主要是为了让每一步的变化能够被观察到。

例如，参考没生效，就沿绑定、编码、adapter 去查；编辑破坏外区，就沿 mask、感受野和最终合成去查；续写逐渐漂移，再看历史分布、窗口和状态记忆。如果一开始同时改这些环节，最后即使效果变了，也很难说明原因。

上面的分解能帮助构造技术方案，却还不能还原 2.5 的具体架构、训练规模或奖励配置。闭源模型可以用不同组合实现相近能力，外部研究给出的是可讨论、可验证的选择。复杂运动和多主体交互也要按任务测试，视觉上逼真，并不等于已经得到了工业仿真的物理真值。

实际接入时，还应单独确认所用入口的时长、延长次数、参考限制和计费。发布时的``多轮延长''与某个入口允许的次数，可以有不同范围；能力展示并不自动等于所有 API、套餐和地域的可用条件。
